\scanpage{51}

\unsure{correspondent} à des syst. d'iso.
\[
\mu_{\mathcal D^*}:f_{\mathcal D^*}\longrightarrow g_{\mathcal D^*},\qquad
\unsure{\mu_{U^*}}\;f_U\longrightarrow g_U.
\]
Cela correspond \uline{conditions} $\tau'=\tau$, $d'_{i'}=d_{i'}$ $\forall i'\in I'$ et la donnée de
\inserted{ceci résulte des autres conditions\unsure{\ldots}}
\margin{\unsure{exprimons} l'\uline{existence} d'un iso. $\mu_{\mathcal D^*}$}

\margin{décrit $\mu_U$}
a) $\boxed{\boxed{\mu\in\mathcal E}}$, définissant un iso entre
\[
f_{\mathcal E}:\mathcal E'\longrightarrow\mathcal E
\quad\text{et}\quad
f'_{\mathcal E}:\mathcal E'\longrightarrow\mathcal E
\]
i.e. \deleted{\unsure{faisant}} tel que
\[
\boxed{f'_{\mathcal E}=\operatorname{int}(\mu)f_{\mathcal E}}\,)
\]

b) $\boxed{\boxed{\mu_{i'}\in\mathcal E_i}}$ $(i'\in I',\ i=\tau(i'))$
\margin{décrit $\mu_{\mathcal D^*}$}
satisfaisant
\[
\boxed{f'_{i'}=\operatorname{int}(\mu_{i'})f_{i'}}\,),
\]
\margin{ceci implique déjà $d_{i'}=d'_{i'}$}
les $\underset{\substack{\unsure{\in}\\\mathcal E}}{\mu},\ \underset{\substack{\unsure{\in}\\\mathcal E_i}}{\mu_{i'}}$ étant liés aux $\alpha,\alpha'$ \deleted{\illegible{1}} $(\unsure{\in\mathcal E|\in\Gamma})$
\inserted{$\alpha\unsure{\in\mathcal E|\in\Gamma}$}
et $\alpha_{i'},\alpha'_{i'}$ $(\in\Gamma)$ de la façon suivante
\[
\boxed{\alpha'=p(\mu)\alpha}
\]
\[
\boxed{\alpha'_{i'}=\rho_i(\mu_{i'})\alpha_{i'}}\qquad\forall i'\in I'
\]
ou encore
\[
\boxed{\begin{cases}
p(\mu)=\alpha'\alpha^{-1}\\
\rho_i(\mu_{i'})=\alpha'_{i'}\alpha_{i'}^{-1}
\end{cases}}
\]
\margin{50}
\scanpage{52}
\mspage{24}

Notons que, comme $\rho_i$ est injectif, les deuxièmes relations, donnant les $\rho_i(\mu_{i'})$, déterminent les $\mu_{i'}$ de \uline{façon unique}, l'existence de ces $\mu_{i'}$ \uline{équivaut aux relations}
\[
\boxed{\alpha'_{i'}\alpha_{i'}^{-1}\in\rho_i(\mathcal E_i)}\qquad i'\in I'.
\]

Quant à la relation $p(\mu)=\alpha'\alpha^{-1}$, elle détermine $\mu$ modulo multiplication \inserted{(à droite disons)} par un \deleted{\illegible{1}} $\mu_0\in\operatorname{Ker}p=\pi$, mais qui doit être tel que l'on ait encore
\[
f'_{\mathcal E}=\operatorname{int}(\deleted{\illegible{1}}\mu\mu_0)f_{\mathcal E}
\]
(en plus de $f'_{\mathcal E}=\operatorname{int}(\mu)f_{\mathcal E}$) ce qui signifie \deleted{exige} que
\[
\operatorname{int}(\mu_0)f_{\mathcal E}=f_{\mathcal E}
\]
i.e.
\[
\boxed{\mu_0\in\operatorname{Centr}_{\unsure{\pi}}\deleted{\illegible{1}}f_{\mathcal E}(\mathcal E')}
\]

C'est donc cela l'indétermination exacte dans le choix d'un $0$-morphisme \unsure{des|les} $1$-morph. \deleted{entre} $f$ et $f'$ de $\mathcal D'$ dans $\mathcal D$.
Si p.ex. on est dans les conditions où
$\pi'\longrightarrow\pi$ induit par $f_{\mathcal E}:\mathcal E'\longrightarrow\mathcal E$ \deleted{est} a une image d'indice fini (image ouverte)
--- cas d'un ``morphisme non constant'' ! ---
alors $\mu_0$ doit centraliser un sous-groupe ouvert de
\margin{51}
\scanpage{53}

$\pi$ --- sauf erreur cela aussi implique $\mu_0=1$, donc il semble bien que \deleted{\illegible{1}} (dans les cas correspondants à des hom dominants de courbes algébriques) \unsure{une|un} \deleted{\unsure{hom}} $0$-morphisme d'un $f$ \unsure{ds} un $f'$ (s'il existe) est unique, i.e. \uline{$\operatorname{Hom}$} $(\mathcal D',\unsure{\mathcal D^*|\mathcal D})$ est une catégorie \uline{discrète}.

On peut interpréter $\mathcal D$ comme $\mathcal E$ muni de sous-groupes $\mathcal E_i\subset\mathcal E$, et de $p:\mathcal E\longrightarrow\unsure{\Gamma_*|\Gamma}$ (noyau $\pi$) enfin d'isomorphismes
\[
k_i:\mathrm T(\overline K)\xrightarrow{\sim}\mathcal E_i\cap\pi=L_i
\]
(commutant à l'action de $\mathcal E_i$).
Si on a un autre syst.
\[
(\mathcal E',\ (\mathcal E'_{i'})_{i'\in I'},\ p':\mathcal E'\longrightarrow\unsure{\Gamma'|\Gamma},\ (k'_{i'})_{i'\in I'}),
\]
\inserted{$\mathcal E'_{i'}\subset\mathcal E'$}
un morphisme du premier dans le second est \uline{défini} par un
\[
f:\mathcal E'\longrightarrow\mathcal E
\]
satisfaisant les conditions suivantes

\inserted{b) $pf$ est conjugué de $\unsure{p'|f'}$ [soit $\alpha\in\Gamma$ tel que $pf=\operatorname{int}(\alpha)p'$]}

a) $\forall i'\in I',\quad\exists i\in I$ tel que $\unsure{f'|f}(\mathcal E'_{i'})$ soit contenu dans un conjugué de $\mathcal E_i$
[soit $\alpha_{i'}\in\mathcal E$ tel que
$f(\mathcal E'_{i'})\unsure{\subset|=}\operatorname{int}(\alpha_{i'}^{-1})(\mathcal E_i)$]
\inserted{i.e. $\operatorname{int}(\alpha_{i'})f(\mathcal E'_{i'})\subset\mathcal E_i$}
\margin{52}
\scanpage{54}
\mspage{25}

\uline{NB} Le $i\in I$ correspondant à $i'\in I'$ est unique, d'où $\tau:I'\longrightarrow I$.
Si les $\alpha_{i'}$ sont choisis, on déduit \deleted{\unsure{$\forall i'$}} hom induits
$f_{i'}:\mathcal E'_{i'}\longrightarrow\mathcal E_i$,
\[
f_{i'}=(\operatorname{int}(\alpha_{i'})\circ f)|\mathcal E'_{i'},
\]
qui induit dès lors, via les $k'_{i'},k_i$, un homomorphisme
\[
f^0_{i'}:\mathrm T(\overline K)\longrightarrow\mathrm T(\overline K).
\]
\margin{\uline{NB} A priori $f^0_{i'}$ est la \unsure{multiplication par} $\chi_{i'}\in\unsure{\widehat{\mathbb Z}^{\,*}|\widehat{\mathbb Z}}$. Le centralisateur dans $\Gamma$ d'un ss-groupe ouvert est réduit à l'unité.}
On exige que $\alpha\in\Gamma$ (qui conjugue $\unsure{f_{\mathcal E}}$ en $\unsure{f'_{\mathcal E}}$, et est \unsure{probablement} uniquement déterminé par cette condition) et $\alpha_{i'}$ (qui est sans doute déterminé mod \deleted{\unsure{conjugaison}} composition à g. avec un élément du normalisateur de $\mathcal E_i$ (du moins du transporteur \inserted{de $\mathcal E_i$} d'un ss-groupe ouvert de $\mathcal E_i$ --- c'est \unsure{itou} si $f$ est iso) puissent être choisis de telle façon que, posant
\[
\beta_{i'}=p(\alpha_{i'})\alpha,
\]
on ait
\[
f^0_{i'}=(d_{i'}\chi(\beta_{i'}))\operatorname{id}_{\mathrm T(\overline K)}
\]
\inserted{$\chi_{i'}=d_{i'}\chi(\beta_{i'})$}
\margin{Cette condition sur le syst. des $\alpha_{i'}$ ne change pas ($\alpha$ restant fixé) lorsqu'on \illegible{3} $\mu_{i'}\alpha_{i'}$, où $\mu_{i'}\in\mathcal E_i$.}
i.e. $d_{i'}\in\mathbb N^*$ est un entier naturel évidemment uniquement déterminé \unsure{\}} quand $\alpha$ et les $\alpha_{i'}$ sont choisis).

\deleted{Remplaçons $\alpha$ par $\alpha'=\alpha_0\alpha$.}
\margin{53}
\scanpage{55}

\unsure{8.} \uline{Réflexion taxonomique.}

Le cas $\dim X=0$ se traduit sur le \illegible{1} du système
$(\mathrm T(\overline K)\xrightarrow{k_i}\mathcal E_i\overset{\rho_i}{\hookrightarrow}\mathcal E\xrightarrow{p}\Gamma)$
\deleted{\illegible{2}} par la condition \deleted{$\unsure{\pi|I}=\varnothing$} $p$ injectif, \deleted{donc se décrit} \deleted{\illegible{1}}
\inserted{et alors $I=\varnothing$ --- donc ce cas se décrit}
simplement par la donnée d'un \uline{sous-groupe} ouvert $\mathcal E\subset\Gamma$, i.e. d'une sous-extension finie $\unsure{L}$ de $\overline K/K$
(NB on aura bien sûr $\unsure{X'|X}=\operatorname{Spec}\unsure{L'|L}$, et les choix \inserted{faits} sur $X$ \inserted{--- i.e. sur $\Pi_{\unsure{\mathcal D^{**}}}\longrightarrow\Pi_U\longrightarrow\Pi_e$ ---} \unsure{aboutissent} à cette description --- \unsure{reviennent|uniquement} ici au choix d'une telle \inserted{$K$-}immersion de $L\deleted{\illegible{1}}=\mathrm H^0(X,\mathcal O_X)$ dans $\overline K/K$)

Le cas $\dim X=1$ se traduit par le fait que $\mathcal E\longrightarrow\Gamma$ n'est pas injectif --- quant à $I$, il peut être vide ou non dans ce cas.
S'il est vide, la description se fait donc simplement en termes d'un hom de groupes profinis $\mathcal E\xrightarrow{p}\Gamma$ (où $\Gamma$ donné d'avance $=\operatorname{Gal}(\overline{\unsure{K|k}}/k)$ \unsure{\}} $=\operatorname{Aut}_{\Pi_e}(\Omega)$).

Supposons donnés maintenant à la fois
\[
\mathcal D=(\mathcal E,p,I,(\rho_i:\mathcal E_i\hookrightarrow\mathcal E),(k_i:\mathrm T(\overline K)\longrightarrow\mathcal E_i))
\]
et
\[
\mathcal D'\unsure{=} (\mathcal E',\unsure{p'},I',(\rho'_{i'}:\ ),(k'_{i'}:\ )),
\]
et revenons à la question de la description des morphismes de $\mathcal D'$ dans $\mathcal D$.
On va distinguer quatre cas, suivant les deux valeurs
\margin{54}
\scanpage{56}
\mspage{26}

possibles $0,1$ de $n=\dim\mathcal D$ et $n'=\dim\mathcal D'$ respectivement.

I) $n=0,n'=0$, alors $I=I'=\varnothing$ et les données se réduisent à des sous-groupes ouverts
$\mathcal E\overset{p}{\hookrightarrow}\Gamma$, $\mathcal E'\overset{p'}{\hookrightarrow}\Gamma$.
Dans la grande description plus haut, \deleted{toutes} les considérations relatives à $I,I'$ tombent et un morphisme $\mathcal D'\longrightarrow\mathcal D$ revient à une classe de \deleted{\unsure{homomorphismes}} couples $(f,\alpha)$, où
$f:\mathcal E'\longrightarrow\mathcal E$ et $\alpha\in\Gamma$, tel que l'on ait
\[
pf=\operatorname{int}(\alpha)\unsure{f|p'}
\]
--- ce qui implique que $f$ est déterminé par $\alpha$, étant induit par $\operatorname{int}(\alpha)$
\inserted{\illegible{2} : $\alpha\in\operatorname{Transp}_{\Gamma}(\mathcal E',\mathcal E)$ \illegible{1} $\operatorname{Transp}_{\Gamma}(L,L')$}
\margin{\unsure{(un)}}
\deleted{\diagram{[Carré barré : $\mathcal E'$ en haut à gauche, $\Gamma$ en bas à gauche et en bas à droite; objet supérieur droit \unsure{$\mathcal E$}; flèche supérieure vers la droite, étiquetée $f$; deux flèches descendantes, étiquetées $p'$ et $p$; flèche inférieure vers la droite.]}}
\diagram{\begin{tikzcd}[ampersand replacement=\&]
\mathcal E' \arrow[r,"f"] \arrow[d,"p'"'] \& \mathcal E \arrow[d,"p"] \\
\Gamma \arrow[r,"\operatorname{int}(\alpha)","\sim"'] \& \Gamma
\end{tikzcd}}
[qui induit donc un hom de $\unsure{L^*|L}\subset\overline K$ dans $L'\subset\overline K$, donc un hom $\operatorname{Spec}(L')\longrightarrow\operatorname{Spec}(L)$ en sens inverse].

La condition que $\alpha,\alpha'\in\Gamma$ définissent le même morphisme $\mathcal D'\longrightarrow\mathcal D$ (le même dans $K$-\unsure{Sch}) se décrit par l'existence d'un $\mu\in\mathcal E$ \deleted{(= stabilisateur de} tel que $\alpha'=\mu\alpha$.
Le fait qu'on donne une corr. $1$-$1$ entre
\[
\operatorname{Hom}_{K}(X',X)\simeq\unsure{\operatorname{Hom}}_{K}(\unsure{L},L')
\]
avec $\operatorname{Hom}(\mathcal D',\mathcal D)$ explicité ainsi, est clair.
\margin{55}
\scanpage{57}

II) $n'=1,n=0$.
\inserted{Ici $n=0$ implique déjà $I=\varnothing$.}
\deleted{Ici} Pour que $\operatorname{Hom}(X',X)\ne\varnothing$, il faut que l'on ait $I'=\varnothing$
[par la condition $f^{-1}(S)=S'$ qui implique $S'=\varnothing$, si $f:(X',S')\longrightarrow(X,S)$]
\inserted{$f^{-1}(\varnothing)=\varnothing$}
--- \deleted{dans} dans la description des morphismes $\mathcal D'\longrightarrow\mathcal D$, cela correspond au fait qu'on ne peut avoir d'application $\tau:I'\longrightarrow I(=\varnothing)$ que si $I'=\varnothing$.
\margin{\unsure{cela on ne peut}}
\deleted{\illegible{1}}
Se bornant donc au cas $I'=\varnothing$ (sinon $\operatorname{Hom}(X',X)=\operatorname{Hom}(\mathcal D',\mathcal D)=\varnothing$ et on est heureux), on trouve donc que $\mathcal D$ revient à la donnée de $\mathcal E\overset{p}{\hookrightarrow}\Gamma$ ss-groupe ouvert, et $\mathcal D'$ à la donnée d'un hom de groupes profinis $p':\mathcal E'\longrightarrow\Gamma$, dont l'image sera un ss-groupe ouvert de $\Gamma$, soit $\Gamma'_0\subset\Gamma$, correspondant à une ss-extension \inserted{finie} $L'$ de $\overline K/K$.
\margin{Se bornant donc au cas}

Les choix faits relatifs à $X'$ impliquent précisément un iso fixé $L'\simeq\mathrm H^0(X',\mathcal O_{X'})$.
Ceci dit, un $K$-morphisme \inserted{de $(X',\varnothing)\longrightarrow(X,\varnothing)$ i.e. de} $X'\longrightarrow X$ revient \inserted{(comme $X$ affine)} à la donnée d'un $K$-morphisme $\operatorname{Spec}L'\longrightarrow X\ (\simeq\operatorname{Spec}L)$ de l'enveloppe affine, i.e. d'un \inserted{$K$-}hom $L\longrightarrow L'$.
(Ce qui ne dépend que de $L,L'$, ou encore \unsure{des} \deleted{\illegible{1}} ss-groupes $\mathcal E\subset\Gamma$ et $\Gamma'_0\subset\Gamma$
\margin{56}
\scanpage{58}
\mspage{27}

de $\Gamma$).
D'autre part la description en termes des diagrammes $\mathcal D,\mathcal D'$ revient à dire qu'un morphisme est déterminé par un couple $(f,\alpha)$, $f:\mathcal E'\longrightarrow\mathcal E$ et $\alpha\in\Gamma$, tels que l'on ait encore commutatif
\diagram{\begin{tikzcd}[ampersand replacement=\&]
\mathcal E' \arrow[r,"f"] \arrow[d,"p'"'] \& \mathcal E \arrow[d,hook,"p"] \\
\Gamma \arrow[r,"\operatorname{int}(\alpha)"] \& \Gamma
\end{tikzcd}}
\[
pf=\operatorname{int}(\alpha)p'
\]
ce qui implique encore que $f$ est déterminé par $\alpha$ (savoir, induit par $\operatorname{int}(\alpha)$), et $\alpha$ étant \unsure{sujet} à la seule condition
\[
\alpha\in\operatorname{Transp}_{\Gamma}(\Gamma'_0,\mathcal E)
\]
--- et $\alpha,\alpha'$ définissent le $=$ hom \unsure{ssi}
\[
\exists\mu\in\mathcal E\quad\text{tel que}\qquad\unsure{\Longleftarrow|\Longleftrightarrow}
\]
\[
\alpha'=\mu\alpha
\]
donc on a encore
\[
\begin{aligned}
\operatorname{Hom}(\mathcal D',\mathcal D)
&\simeq\mathcal E\backslash\operatorname{Transp}_{\unsure{\Gamma|\mathcal E}}(\Gamma'_0,\unsure{\Gamma|\mathcal E})\deleted{\illegible{1}}\\
&\simeq\operatorname{Transp}_{\Gamma}(L,L')/\operatorname{Fix}_{\Gamma}(L)\\
&\simeq\operatorname{Hom}_{\unsure{K\text{-}\mathrm{alg}|K\text{-}\mathrm{cor}}}(L,L')
\end{aligned}
\]
donc on trouve encore dans ce cas la pleine fidélité.

En fait, on aurait pu traiter ensemble
\margin{57}
\scanpage{59}

les deux cas $(n'=0,n'=1)$ où $n=0$, et la description des $\operatorname{Hom}(X,X')$ \inserted{en termes de $\mathcal D,\mathcal D'$} est valable d'ailleurs, sans \deleted{supposer que} faire des hypothèses draconiennes (projectif, lisse, $\dim\leq1$, anabélien) sur $X'$.
\deleted{Il suffit que $X'$ soit un schéma connexe quasi-compact quasi-séparé sur $K$, alors les morphismes $X'\longrightarrow X=\operatorname{Spec}(L)$ correspondent aux morphismes $\operatorname{Spec}A'\longrightarrow X$ de l'enveloppe affine de $X'$ en $X$, définie comme Spectre de $A'=\mathrm H^0(X',\mathcal O_{X'})$, i.e. aux \unsure{$K$-hom} de $K$-algèbres
\[
L\longrightarrow A'.
\]}

En effet, comme $X$ est \inserted{fini} étale sur $K$, l'image inverse \inserted{$X_{X'}=$} \deleted{\illegible{1}} $X'\times_K\unsure{X}$ sur $X'$ est \inserted{fini} étale sur $X'$, et
\[
\operatorname{Hom}_K(X,X')\simeq\deleted{\operatorname{Hom}\illegible{1}}\operatorname{Sections}(X_{X'}/X')
\]
peut se décrire en termes de la catégorie $\unsure{\widehat{\Pi_1(X')}|\widehat{\Pi}_1(X')}$ $(=\Pi_U)$ des \unsure{syst.|revêt.} locaux finis sur $X'$ i.e. des syst. locaux finis sur $\Pi_U$.
Suivant cette description, on aboutit à la
\margin{58}
\scanpage{60}
\mspage{28}

description commune donnée dans I, II, sans aucune hyp. autre sur $X$ que \deleted{la connexité} la $0$-connexité (pour pouvoir décrire $\Pi_U=\Pi_{X'}$ sur $\Pi_e$ en termes d'un hom de groupes $p':\mathcal E'\longrightarrow\Gamma\ldots$).

Quand on se donne $(X',S')$, avec $S'$ pas nécessairement vide, $X'$ $0$-connexe (et \deleted{dim} $X'$ lisse de dim $1$ aux pts de $S'$ (et il suffirait $=$ que $S'$ \unsure{ne} \unsure{disconnecte|déconnecte} pas localement \inserted{(ét)} $X'$, $\ldots$),
\inserted{donc $X'-S'$ $0$-connexe}
alors la \deleted{partie} donnée
$\Pi_1(X'-S')\longrightarrow\Pi_1(e_K)$ est décrite encore par $\mathcal E'\xrightarrow{p'}\Gamma$ (indépendamment de la considération des $\mathcal E'_i$ etc) et par suite $\operatorname{Hom}_K(X'-S',X)$ est décrit comme on vient de dire.
Mais on a (avec les hypothèses de normalité de $X'$ aux pts de $S'$)
\[
\operatorname{Hom}_K(X'-S',X)=\operatorname{Hom}_K(X',X).
\]

Ceci encourage à modifier la définition des morphismes $(X',S')\longrightarrow(X,S)$ donnée \unsure{au|par} début, via $f:X'\longrightarrow X$ avec $f^{-1}(S)_{\mathrm{r\acute ed}}=S'$, en exigeant ceci \inserted{non pour $f$, mais} \unsure{seulement} pour la restriction de $f$ \unsure{à} la réunion des comp. irréd. de $X'$
\margin{59}
